\section{Discusión}\label{sec:discusion}

Los resultados definitivos fuera de muestra permiten contrastar las decisiones de diseño de SM-CFHR frente a la evidencia empírica y delimitar con precisión dónde aporta ventaja real y dónde falla frente a alternativas locales o no lineales.

\subsection{Interpretación empírica por régimen, horizonte y estrategia}

\begin{enumerate}
    \item \textbf{Ventaja concentrada en transferencia espaciotemporal con meteorología:} la combinación de normalización causal por $24\,000\,C_{i,t}$, regularización anisotrópica y cotas monótonas logra rango $1{,}0/23$ y $5/5$ victorias Holm bajo el estimando $\mathrm{MAE}_{\mathrm{macro}}$ en \texttt{D+1 ET} y \texttt{RecGEFSv12 ET} ($7$--$28$~días). Al evaluar plantas nunca vistas durante el ajuste, los regresores entrenados en energía bruta ($\text{kWh/día}$) sufren la distorsión de escala de $39{,}6\times$, mientras SM-CFHR transfiere relaciones adimensionales consistentes entre zonas geográficas.
    \item \textbf{Limitaciones en intraplanta y caída en estrategia directa sin clima:} cuando existe historial propio de la planta (\texttt{IP}), las bibliotecas locales o flexibles (\texttt{BayesianRidge}, \texttt{k-NN}, \texttt{Random Forest}) aprovechan relaciones no lineales y niveles específicos de energía sin necesitar transferencia de escala. En particular, en \texttt{DirectoSinClima\_IP} ($7$--$28$~días), el ajuste lineal adimensional sobre acumulados sin señal NWP se degrada hasta $+24{,}24\%$ frente a \texttt{k-NN}, constituyendo un resultado negativo claro del método.
    \item \textbf{Compromiso entre configuración fija y HPO:} \texttt{SM-CFHR-Fixed} ($B=0$) supera a \texttt{SM-CFHR-HPO} ($B=25$) en $19/28$ escenarios ($14/14$ en ET; $4$ victorias Holm frente a $1$ derrota) con un coste de pared $59\times$ menor. La varianza de selección en $3$ pliegues internos penaliza a \texttt{HPO} especialmente en transferencia espacial y horizontes largos, mostrando que el sesgo inductivo fijo actúa como regularizador efectivo.
    \item \textbf{Sensibilidad de escala ($\text{kWh}$ vs.\ CF) y subcobertura conformal:} el contraste entre rango medio $2{,}93/23$ en $\text{kWh}$ y $4{,}00/23$ en $\mathrm{MAE}_{\mathrm{CF,macro}}$ evidencia que las centrales de mayor potencia influyen fuertemente en las métricas en energía absoluta. Asimismo, las coberturas conformes empíricas ($83{,}24\%$ para $90\%$ nominal; $69{,}25\%$ para $80\%$) confirman que la no estacionariedad y deriva entre bloques violan la intercambiabilidad estándar \cite{barber2023nonexchangeable}.
\end{enumerate}

\subsection{Implicaciones estructurales de las decisiones de diseño}

La arquitectura de SM-CFHR introduce cuatro sesgos inductivos explícitos: normalización causal por $24\,000\,C_{i,t}$, escalado MAD con pérdida de Huber, penalización $L_2$ anisotrópica ($\alpha_{\mathrm{lin}}, \alpha_{\mathrm{meteo}}$) y restricciones de signo no negativo sobre la persistencia reciente y la irradiación pronosticada. Esta elección privilegia la estabilidad paramétrica, la convexidad del ajuste de un paso y la invariancia de escala entre centrales de $9{,}8$ a $388{,}0$~MW, a costa de restringir la flexibilidad funcional frente a relaciones no lineales complejas entre variables meteorológicas. Asimismo, en el despliegue recursivo multihorizonte ($H \in \{7,14,21,28\}$), la realimentación de predicciones propias mediante ventana deslizante y amortiguamiento $\gamma_h = 0{,}85^{h-1}$ transforma un problema interno convexo en una composición no lineal de horizonte $H$, haciendo indispensable que la validación interna de hiperparámetros reproduzca exactamente esa misma dinámica recursiva.

En el plano experimental, cuatro aspectos condicionan la interpretación de cualquier resultado futuro: (1) la asimetría de presupuestos de búsqueda entre \texttt{SM-CFHR-Fixed} ($B=0$), los modelos de referencia (rejilla de $4$ configuraciones) y \texttt{SM-CFHR-HPO} ($B=25$ ensayos TPE), razón por la cual la comparación primaria se fija en $B=0$; (2) la reutilización adaptativa del panel histórico 2024--2026 durante el desarrollo de los modelos previos y la consolidación de SM-CFHR \cite{cawley2010overfit}; (3) el sesgo potencial de selección posterior al comparar en pruebas pareadas frente al mejor competidor de biblioteca identificado sobre la propia muestra de evaluación, que exige complementar dicho contraste con la línea base preespecificada \texttt{MM30}, comparaciones exploratorias frente a referencias individuales de biblioteca y pruebas múltiples (SPA y MCS con interpretación estricta), junto con la distinción entre la métrica física exportada en $\text{kWh/día}$ (ponderada por capacidad instalada) y métricas adimensionales en factor de capacidad ($\mathrm{MAE}_{\mathrm{CF}}$, calculable post-hoc sobre las predicciones OOF); y (4) el tamaño muestral efectivo reducido en horizontes de $21$ y $28$ días con orígenes no solapados, donde la autocorrelación residual limita la potencia de los contrastes pareados ($T_{\mathrm{eff}} < 2H$).

\subsection{Limitaciones de generalización}\label{sec:limitaciones_gen}

El estudio utiliza exclusivamente $17$ centrales fotovoltaicas de un mismo dominio geográfico y regulatorio (el Sistema Interconectado Nacional de Colombia) y no evalúa conjuntos de datos externos independientes de naturaleza distinta. Esta limitación científica exige distinguir con rigor cuatro dimensiones de generalización, ninguna de las cuales implica automáticamente a las demás:

\begin{enumerate}
    \item \textbf{Generalización espacial:} mide el desempeño sobre centrales o zonas geográficas apartadas del entrenamiento dentro de la misma cohorte colombiana mediante validación \textit{leave-one-group-out}. Aunque este protocolo evita la fuga entre plantas vecinas situadas a menos de $100\text{ km}$, la flota comprende solo $G=7$ conglomerados espaciales \cite{cameron2015cluster}, por lo que la extrapolación a regímenes climáticos distintos (como zonas desérticas extratropicales o latitudes altas con marcada estacionalidad solar) no queda demostrada.
    \item \textbf{Generalización temporal:} mide el desempeño sobre periodos no utilizados para ajustar parámetros dentro de la ventana enero de 2024--octubre de 2026. Dado que la arquitectura fue fijada tras inspeccionar la estructura exploratoria de ese mismo intervalo histórico, la evaluación OOF mide generalización temporal retrospectiva bajo particiones causales, no desempeño prospectivo sobre años futuros ni cambios estructurales de red posteriores a octubre de 2026.
    \item \textbf{Generalización entre datasets:} mide el desempeño frente a fuentes de medición, resoluciones temporales, sistemas de telemetría o procesos de generación de datos diferentes (por ejemplo, datos SCADA a nivel de inversor, otras redes eléctricas o distintas mallas NWP). Al no haberse evaluado ningún conjunto de datos externo independiente de XM/Open-Meteo/GEFSv12, no puede afirmarse invarianza frente a cambios de fuente o protocolo de registro.
    \item \textbf{Generalización entre dominios:} concierne a la transferibilidad de los principios de SM-CFHR fuera del pronóstico fotovoltaico. La normalización por $24\,000\,C_{i,t}$, el acotamiento físico al intervalo $[0,1]$, el amortiguamiento geométrico de persistencia y las restricciones de no negatividad sobre la persistencia de $30$ días y la irradiación pronosticada dependen específicamente de la física de generación solar. Su adaptación a otros dominios de regresión requeriría redefinir justificadamente las transformaciones de escala, los grupos de regularización y las restricciones de signo pertinentes; por ello, la evaluación diseñada no demuestra universalidad del estimador.
\end{enumerate}
